% TOP_PROBLEM: {"id": 1431, "title": "Characterization of Word-Representable Planar Graphs", "category_id": 3, "category": {"id": 3, "name": "graph_theory", "display_name": "Graph Theory", "description": "Problems involving graphs, networks, and their properties.", "slug": "graph-theory", "order_index": 3, "created_at": "2026-07-31T15:26:25.671Z"}, "status": "open"}
% FIRST_POSTED: null
% ENTRY_KIND: statement_only
% Imported from: batch07.tex; UnsolvedMath ID 1431
\documentclass[11pt]{article}
\usepackage[margin=25mm]{geometry}
\usepackage{fontspec}
\setmainfont{Latin Modern Roman}
\usepackage{amsmath,amssymb,mathtools}
\usepackage{unicode-math}
\setmathfont{Latin Modern Math}
\usepackage{xurl,hyperref}
\hypersetup{colorlinks=true,urlcolor=blue}
\setcounter{secnumdepth}{0}
\setlength{\parindent}{0pt}
\setlength{\parskip}{5pt}
\setlength{\emergencystretch}{3em}
\newcommand{\sref}[2]{\hyperlink{source:#1:#2}{[S#2]}}
\newcommand{\cataloguescope}[1]{\paragraph{Recorded target.}#1}
\begin{document}
% UnsolvedMath ID 1431; source catalogue code GRAPH-044
\setcounter{section}{1430}
\section{Characterization of Word-Representable Planar Graphs}\label{1431}
{\small\sffamily UnsolvedMath ID 1431 · Graph Theory}
\subsection{Definitions and mathematical statement}

\paragraph{Shared notation.}$\mathbb N=\{1,2,\ldots\}$, and $\mathbb Z,\mathbb Q,\mathbb R,\mathbb C$ denote the integers, rationals, reals and complex numbers. Any other conventions are specified locally.


Let $G=(V,E)$ be a finite simple undirected graph. A word over $V$ is a finite sequence of vertices in which every vertex occurs. For distinct $x,y\in V$, let $w|_{\{x,y\}}$ be the word obtained by deleting every letter other than $x,y$. They alternate when this restriction has no equal consecutive letters, that is, it has the form $xyxy\cdots$ or $yxyx\cdots$, of either even or odd length. The word $w$ represents $G$ when, for every distinct $x,y$,
\[
 xy\in E\quad\Longleftrightarrow\quad x,y\text{ alternate in }w.
\]
Write $\mathrm{WR}(G)$ when such a word exists.

A graph is planar if it can be drawn in $\mathbb R^2$ with different edges intersecting only at a shared endpoint.
\paragraph{Statement recorded in the supplied source.}
Give a complete structural characterization, up to graph isomorphism, of
\[
 \{G:G\text{ is finite and planar},\ \mathrm{WR}(G)\}.
\]
The criterion must apply to arbitrary planar graphs, not only triangle-free graphs or near-triangulations. The general characterization by existence of a semi-transitive orientation is already known; the recorded request is for a structural description of the planar class. \sref{1431}{2}
\subsection{Short English statement}
Describe exactly which finite planar graphs can be encoded by a word whose alternating pairs are precisely their edges.
\subsection{Sources}
\begin{description}
\item[\hypertarget{source:1431:1}{S1}] ULAM.AI and the UnsolvedMath Contributors, UnsolvedMath: A Curated Collection of Open Mathematics Problems, record 1431 (GRAPH-044). CC BY 4.0. \url{https://huggingface.co/datasets/ulamai/UnsolvedMath}.
\item[\hypertarget{source:1431:2}{S2}] Suchanda Roy and Ramesh Hariharasubramanian, Characterization of Word-Representable Near-Triangulations, arXiv:2605.25733v1 (2026), Sections 1--2, near-triangulation definition and word-representability. \url{https://arxiv.org/pdf/2605.25733v1}.
\item[\hypertarget{source:1431:3}{S3}] Magnus M. Halldórsson, Sergey Kitaev and Artem Pyatkin, Semi-Transitive Orientations and Word-Representable Graphs, arXiv:1501.07108v1 (2015), Section 1, definitions and general orientation characterization. \url{https://arxiv.org/pdf/1501.07108v1}.
\end{description}
\label{end:1431}
\end{document}
